\section{Deterministic comparison fields}\label{m:sec:barriers}

We now construct positive fields that can be compared across all the
local scales.  The construction uses only the nuclear measure.  Its
output is a family of continuous offsets satisfying a nonlinear
subsolution inequality, together with a strictly positive gap when the
scale doubles.  The gap will allow the quantum field to replace a
comparison field in the next stage of the proof.

Fix a nuclear system and a small parameter $r_*>0$.  Its value in terms
of the excess charge, together with the dyadic scales used for the
observations, will be chosen in Section~\ref{m:sec:conditioning}.
Every smallness threshold for $r_*$ in the present deterministic
construction is universal and independent of the nuclear system.

Throughout this section, $p=3/2$, $k=((5/3)\cT)^{-3/2}$, and
$\cF=4\pi k$.  Recall that $K(x)=|x|^{-1}$ and that $a_x$ is the nuclear
scale of Lemma~\ref{m:lem:screen-scale}.  All constants in this section are independent of the nuclear
configuration and charges.  They may depend on the fixed localization
constants and the smooth radial function $g$ used for wave packets.

\subsection{Global Thomas--Fermi fields}

The variational and comparison framework for point and smoothed
sources is classical
\citep[Theorems~II.14, II.17--II.18, and IV.7--IV.10]{LiebSimon1977}.
We give a proof on the Coulomb-space cone used here, including the decay
needed for comparison on an unbounded domain.  If $\zeta$ is a finite
sum of positive point charges and nonnegative smooth compactly supported
radial charges, write $V_\zeta=K*\zeta$.  A radial charge is specified by
its center, profile, and total mass.  Set
\[
 \mathcal X_+=\{\rho\ge0:\rho\in L^{5/3}(\mathbb R^3)
                  \cap(\dot H^1(\mathbb R^3))^*\}.
\]
By Lemma~\ref{lem:loc-coulomb}, the Coulomb energy of an element of this
cone is the squared norm
$D(\rho)=\|\nabla K*\rho\|_2^2/(8\pi)$.

\begin{lemma}[Global comparison problem]\label{m:lem:barrier-global-tf}
Let $\zeta$ be as above, with total mass $Q>0$.  The functional
\[
 \mathcal T_\zeta(\rho)
   =\cT\int\rho^{5/3}-\int V_\zeta\rho+D(\rho),
       \qquad \rho\in\mathcal X_+,
\]
has a unique minimizer.  Its field $W=V_\zeta-K*\rho$ satisfies
\begin{equation}\label{m:eq:global-tf-equation}
 \rho=kW_+^p,\qquad
 \Delta W=-4\pi\zeta+\cF W_+^p,\qquad
 W\ge0,\qquad \int\rho\le2Q.
\end{equation}
The screening potential $K*\rho$ is continuous and tends to zero at
infinity.  Consequently $W$ tends to zero at infinity and is continuous
away from point charges.
\end{lemma}

\begin{proof}
Smear the point charges at any fixed positive widths and leave the
smooth charges unchanged.  This decomposes $V_\zeta=V_s+V_c$ with
$V_s\in\dot H^1$ and compactly supported $V_c\in L^{5/2}$.  Thus
\[
 \left|\int V_s\rho\right|
     \le C\|\nabla V_s\|_2\sqrt{D(\rho)},\qquad
 \left|\int V_c\rho\right|
     \le\|V_c\|_{5/2}\|\rho\|_{5/3}.
\]
For the fixed source and smearing widths, both displayed potential norms
are finite. Young's inequality proves coercivity in both spaces defining
$\mathcal X_+$.  A bounded minimizing sequence has weakly convergent
subsequences in these two reflexive spaces.  Their limits agree as
distributions, since both converge against compact smooth tests;
nonnegativity passes to this limit.  The two linear terms are weakly
continuous in their respective spaces, and both positive terms are
weakly lower semicontinuous.  This proves existence.  The strictly
convex kinetic term proves uniqueness.

Nonnegative smooth compact variations give
\[
 d:=\tfrac53\cT\rho^{2/3}-V_\zeta+K*\rho\ge0
 \quad\text{almost everywhere}.
\]
For bounded smooth compactly supported $\varphi$, the variations
$(1+t\varphi)\rho$ are admissible for sufficiently small positive and
negative $t$.  Indeed $|\varphi\rho|\le\|\varphi\|_\infty\rho$, and
positivity of the Coulomb kernel bounds its Coulomb energy.  The
nonnegative pairing property in Lemma~\ref{lem:loc-coulomb} therefore
justifies differentiating these variations and gives
$\int d\varphi\rho=0$.  Hence $d\rho=0$, proving the Euler equation in
\eqref{m:eq:global-tf-equation} and its distributional Laplacian.

To establish finite mass without assuming it, let the source components
have centers $R_m$ and masses $Q_m$, and put
\[
 U(x)=\sum_m\frac{Q_m}{|x-R_m|},\qquad w(x)=U(x)^{-1}.
\]
Newton's theorem gives $V_\zeta\le U$.  Where $\rho>0$, the Euler
equation gives $K*\rho\le V_\zeta$.  For any ball $B$, multiply this
inequality by $1_Bw\rho$ and restrict the screening integral below to
$B$.  Write
$q_m(x)=w(x)Q_m/|x-R_m|$, so that $\sum_mq_m=1$.  The triangle
inequality gives
\[
 \frac{w(x)+w(y)}{|x-y|}
     \ge\sum_m\frac{q_m(x)q_m(y)}{Q_m}.
\]
The resulting symmetrization and Cauchy--Schwarz yield, with
$m_B=\int_B\rho$,
\[
 m_B\ge\frac12\sum_m\frac{(\int_Bq_m\rho)^2}{Q_m}
       \ge\frac{m_B^2}{2Q}.
\]
All integrals in this calculation are local and finite; in particular
$w$ is bounded on $B$, and the local Coulomb energy is finite.  Increasing
$B$ proves $\int\rho\le2Q$.

For completeness, $L^1\cap L^{5/3}$ implies continuity and decay of the
screening potential.  The kernel restricted to a fixed ball belongs to
$L^{5/2}$, so its pairing with translates of $\rho$ is continuous.  The
complementary kernel contributes at most $\|\rho\|_1/R$ when the
cutoff radius is $R$.  For fixed $R$, the local $L^{5/3}$ norm of $\rho$
on a ball whose center tends to infinity tends to zero.  First letting
the center tend to infinity and then $R\to\infty$ proves decay.

Finally, $(W+\varepsilon)_-$ is compactly supported away from point
charges for every $\varepsilon>0$: $W\to0$ at infinity, and its positive
Coulomb singularities dominate the continuous screening potential.
On the support of this test $\rho=0$, and hence $\Delta W\le0$.
Testing the equation gives
$\int|\nabla(W+\varepsilon)_-|^2\le0$.
The test is justified by compact $H^1$ approximation away from the
singularities.  Letting $\varepsilon\downarrow0$ proves $W\ge0$.
\end{proof}

\begin{lemma}[Point fields and mixtures]\label{m:lem:barrier-mixture}
For one point charge $Q>0$ at $R$, the field in
Lemma~\ref{m:lem:barrier-global-tf} satisfies
\begin{equation}\label{m:eq:single-field-lower}
 W(x)\ge c\min\left(\frac{Q}{|x-R|},|x-R|^{-4}\right).
\end{equation}
For a single point charge or radial smooth charge of mass $Q$, its
screening potential is bounded by $CQ^{4/3}$.

More generally, select source components of masses $Q_i$, truncate them
to masses $q_i\le Q_i$, and put $q=\sum_iq_i>0$.  Let $W_i$ be the
single-source field of mass $q$ at the $i$th center with that source's
normalized profile.  If $W$ is the full-source field, then
\begin{equation}\label{m:eq:tf-mixture}
 \sum_i\frac{q_i}{q}W_i\le W.
\end{equation}
\end{lemma}

\begin{proof}
For a single source, Newton's theorem and the Euler equation give
$\rho(z)\le k(Q/|z-R|)^{3/2}$.  For every $x,R$ and $s>0$,
\begin{equation}\label{m:eq:local-power-potential}
 \int_{B(x,s)}\frac{dz}{|x-z|\,|z-R|^{3/2}}
       \le Cs^{1/2}.
\end{equation}
To see the uniformity in $R$, apply H\"older with exponent $q_1\in(3/2,2)$
to the second factor and its conjugate $q_1'\in(2,3)$ to the first.
The integrals of both radial powers over $B(x,s)$ are bounded by their
integrals over centered balls of the same radius.  Their powers of $s$
combine to $s^{1/2}$.  Split $K*\rho(x)$ at $s=Q^{-1/3}$.  The near
part is at most $CQ^{3/2}s^{1/2}$ by
\eqref{m:eq:local-power-potential}, and the far part at most $2Q/s$ by
Lemma~\ref{m:lem:barrier-global-tf}.  Both are $CQ^{4/3}$.

For a point charge this gives $W\ge Q/(2d)$ for
$d=|x-R|\le cQ^{-1/3}$.  Outside that sphere set $v=\lambda d^{-4}$.
For sufficiently small universal $\lambda>0$,
\[
 \Delta v=12\lambda d^{-6}\ge \cF v^{3/2},
\]
and its boundary value is below the preceding lower bound.  On the
exterior domain, testing the equation for $v-W$ with
$(v-W-\varepsilon)_+$ gives a nonpositive gradient square bounded
below by zero.  This test has compact support since both fields decay.
Thus $v\le W$.  Together with the inner estimate this proves
\eqref{m:eq:single-field-lower}.

For the mixture $F=\sum_i(q_i/q)W_i$, convexity gives
\[
 \Delta F\ge-4\pi\zeta_{\rm truncated}+\cF F^p
            \ge-4\pi\zeta+\cF F^p.
\]
Compare with the equation for $W$ and test with
$(F-W-\varepsilon)_+$.  The support is compact.  Near a point charge
whose mass was strictly decreased, $F-W$ tends to $-\infty$, so the
test vanishes there.  If the masses agree, the singular terms cancel
and the difference is locally $H^1$.  The remaining density terms are
locally in $L^{5/3}$, so compact $H^1$ approximation is valid.  Monotonicity
of the power gives $F\le W$.
\end{proof}

\subsection{Fields adapted to the nuclear scale}

Fix a sufficiently small constant $c_1\in(0,1/100)$.  Smear the $m$th
nucleus with the radial probability $g_{c_1a_{R_m}}^2$ and denote the
resulting density by $\nu_s$.  Put
\begin{equation}\label{m:eq:nuclear-splitting}
 V_s=K*\nu_s,\qquad V_c=V-V_s\ge0.
\end{equation}
The subscript $c$ records the part confined near the nuclei.  A nucleus
contributing to $V_c(x)$ lies within $2c_1a_x$ of $x$.  Its scale is
comparable to $a_x$, and the total contributing charge is at most
$a_x^{-3}$.  Therefore
\begin{equation}\label{m:eq:smoothed-source-bounds}
 \nu_s(x)\le Ca_x^{-6},\qquad V_c\in L^{5/2}_c(\mathbb R^3).
\end{equation}
The smooth potentials of nuclei within any of the fixed local multiples
of $a_x$ used below sum to at most $Ca_x^{-4}$.  Indeed each such width
is comparable to $a_x$, each unit-mass smooth potential is at most
$C/a_x$, and the total local charge is at most $a_x^{-3}$.

Let $W_0,\rho_0$ be the field and density for the original nuclear measure
$\nu$.  Its continuous offset is
\[
 h_0=W_0-V_c=V_s-K*\rho_0.
\]
For a fixed positive $\theta$, to be chosen below, let $W_s$ be the field
for the smooth source $\theta\nu_s$.

\begin{lemma}[Two reference fields]\label{m:lem:barrier-fields}
There are fixed constants $c_0>0$, $\theta>0$, $c_\theta>0$, and $C<\infty$
such that, off the nuclei,
\begin{equation}\label{m:eq:reference-field-bounds}
 W_0(x)\ge c_0a_x^{-4},\qquad |h_0(x)|\le Ca_x^{-4},
\end{equation}
and everywhere
\begin{equation}\label{m:eq:smooth-reference-bounds}
 c_\theta a_x^{-4}\le W_s(x)\le Ca_x^{-4},\qquad
 |\nabla W_s(x)|\le Ca_x^{-5},\qquad W_s(x)\le1.01W_0(x).
\end{equation}
Here $h_0$ is continuous, and $W_s$ is locally $C^1$.  Both fields decay
at infinity.
\end{lemma}

\begin{proof}
By the scale definition,
$\nu(B(x,2La_x))>(2a_x)^{-3}$.  Truncate charges in this ball to total
$q=(2a_x)^{-3}$ and apply Lemma~\ref{m:lem:barrier-mixture}.  Their distances
from $x$ are at most $2La_x$, so
\eqref{m:eq:single-field-lower} gives the first bound in
\eqref{m:eq:reference-field-bounds} with a fixed $c_0>0$.

If some charges contribute to $V_c(x)$, denote their total by
$q\le a_x^{-3}$.  The screening estimate for each single field of
charge $q$, followed by the mixture comparison, gives
\[
 W_0(x)\ge\sum_{\rm contributors}\frac{z_m}{|x-R_m|}-Cq^{4/3}
          \ge V_c(x)-Ca_x^{-4}.
\]
If there are no contributors, the same conclusion follows from
$W_0\ge0$.

For the upper bound, set $a=a_x$, $s=4a$, and let $V_{\rm loc}$ be the
potential of nuclei within $20a$ of $x$.  On $B(x,s)$ the function
\[
 V_{\rm loc}(z)+\frac{C s^4}{(s^2-|z-x|^2)^4}
\]
is a supersolution of the field equation when $C$ is large: the
Laplacian of the second term is at most
$80Cs^6/(s^2-|z-x|^2)^6$, whereas its $p$th power has the same
spatial dependence and coefficient $C^{3/2}$.  The boundary blowup
dominates $W_0-V_{\rm loc}$.  The positive-difference test is legitimate
because the point singularities cancel.  Thus
$W_0(x)\le V_{\rm loc}(x)+Ca^{-4}$.  Subtracting $V_c(x)$ and using
the smooth local potential bound proves the upper offset estimate.
The expression defining $h_0$ proves its continuity.

For $W_s$, the same bump without $V_{\rm loc}$ is a supersolution after
its coefficient is enlarged to absorb the locally bounded source
$4\pi\theta\nu_s$.  This proves $W_s\le Ca^{-4}$.  Its Laplacian is
bounded by $Ca^{-6}$ on a local ball.  Subtract the Newton potential
of this source restricted to that ball.  Its gradient is bounded by
$Ca^{-5}$, since the integral of $|z|^{-2}$ on a ball of radius $a$ is
$Ca$.  The remaining harmonic function is bounded by $Ca^{-4}$ on a
smaller ball; its interior gradient bound proves the stated estimate
and local $C^1$ regularity.

For the lower bound, select and truncate the same charges as in the
first paragraph, before multiplying them by $\theta$.  Their smearing
widths are at most $3c_1a$, because their centers are within $2La$.
A single such blob of charge $\theta(2a)^{-3}$ has bare potential at
$x$ at least $c_L\theta a^{-4}$: every source point is within
$(2L+3c_1)a$ of $x$.  Its screening potential is at most
$C\theta^{4/3}a^{-4}$.  Choose $\theta$ sufficiently small and apply
the mixture comparison to obtain $W_s\ge c_\theta a^{-4}$.

Decrease $\theta$ further if necessary.  The source bound and the lower
bound for $W_0$ imply
\[
 4\pi\theta\nu_s
   \le \cF\bigl(1.01^{3/2}-1.01\bigr)W_0^{3/2}.
\]
Consequently $1.01W_0$ is a supersolution for the field of
$\theta\nu_s$; its extra negative point sources have the required
sign.  The positive-difference test for $W_s-1.01W_0$, supported away
from those point singularities, proves the final comparison.
Decay follows from Lemma~\ref{m:lem:barrier-global-tf}.
\end{proof}

The reference fields now have positive lower bounds at every scale.
To compare them with the local quantum constructions, we also need to
control variation over a much smaller ball.  The offset is the right
quantity for this purpose: subtracting $V_c$ removes every nuclear
singularity from the function being compared.

\begin{lemma}[Oscillation of offsets]\label{m:lem:barrier-oscillation}
Fix $y$ and write $a=a_y$.  Let $h=h_0$, or let $h=W-V_c$ for an interior
Thomas--Fermi patch centered at $y$ satisfying the hypotheses of
Lemma~\ref{m:lem:screen-patch}.  Then $h$ has a continuous representative.  For
$0<s\le a/10$ and $|z-y|\le s$,
\begin{equation}\label{m:eq:offset-oscillation}
 |h(z)-h(y)|\le Ca^{-4}(s/a)^{1/2}+C(s/a)h(y)_-.
\end{equation}
For $h=h_0$ the final term can be omitted.
\end{lemma}

\begin{proof}
In either case the local upper field estimate in
\eqref{m:eq:screen-interior-upper} or Lemma~\ref{m:lem:barrier-fields}
gives $h\le Ca^{-4}$
on a fixed ball about $y$.  The patch core is absent there.  The Euler
equation, the bound on $\nu_s$, and the weighted power inequality give
\begin{equation}\label{m:eq:offset-laplacian-bound}
 |\Delta h(z)|\le Ca^{-6}
     +Ca^{-3/2}\sum_{\rm loc}z_m|z-R_m|^{-3/2}.
\end{equation}
Here all centers in the sum belong to one fixed local ball and their
total charge is at most $a^{-3}$.  The same estimate holds for $h_0$
by Lemma~\ref{m:lem:barrier-fields}.

Take the Newton particular solution for this Laplacian restricted to
a ball of radius $a$.  Equation~\eqref{m:eq:local-power-potential} bounds
it by $Ca^{-4}$ on a smaller ball.  Its variation over distance $s$
is at most $Ca^{-4}(s/a)^{1/2}$.  In detail, the part within distance
$Cs$ of either evaluation point is at most
\[
 C\bigl(a^{-6}s^2+a^{-9/2}s^{1/2}\bigr).
\]
On a dyadic ring of radius $l\ge Cs$, the difference of the two Coulomb
kernels is at most $Cs/l^2$, and the source mass in that ring is at most
$C(a^{-6}l^3+a^{-9/2}l^{3/2})$.  Summing the rings up to radius $a$
gives the asserted variation.

The harmonic remainder is bounded above by $Ca^{-4}$.  Harnack's
inequality and the harmonic gradient estimate, applied to this upper
constant minus the remainder, bound its variation by
$C(s/a)(a^{-4}+h(y)_-)$.  Combining the estimates gives
\eqref{m:eq:offset-oscillation}, as well as continuity.  For $h_0$,
Lemma~\ref{m:lem:barrier-fields} bounds its negative part by $Ca^{-4}$,
which is absorbed by the first term.
\end{proof}

\subsection{Averaging at a smaller radius}

Fix $w=10^{-5}$, and use the small parameter $r_*$ fixed above.
Define the smearing widths, center kernels,
and averaging operator by
\begin{equation}\label{m:eq:master-kernel}
 t_z=c_1a_z\min(a_z,r_*)^w,\qquad
 \mathcal K_z(y)=g_{t_z}^2(y-z),\qquad
 Rf(y)=\int\mathcal K_z(y)f(z)\,dz,
 \qquad \tau_y=t_y/a_y.
\end{equation}
Thus $R$ smooths a mass at its own center-dependent width.  It preserves
total mass, although it need not preserve the constant function exactly.

\begin{lemma}[Kernel and potential estimates]\label{m:lem:barrier-kernel}
Uniformly as $r_*\downarrow0$, a center contributing to $Rf(y)$ satisfies
$|z-y|\le2t_y$ and $t_z/t_y=1+O(r_*^w)$.  Moreover,
\begin{equation}\label{m:eq:master-kernel-bounds}
 R1(y)=1+O(r_*^w),\qquad
 \|\nabla_z\mathcal K_z(y)\|_\infty\le Ct_y^{-4},\qquad
 \|\nabla_z\mathcal K_z(y)\|_2\le Ct_y^{-5/2}.
\end{equation}
The norms in these gradient estimates are taken in the center variable $z$.
The operator $R$ is bounded on $L^{5/3}$ and preserves integrals.
The correction
\begin{equation}\label{m:eq:potential-correction}
 q_0=K*(\rho_0-R\rho_0)
\end{equation}
is nonnegative, continuous, belongs to $\dot H^1$, and tends to zero at
infinity.  It satisfies
\begin{equation}\label{m:eq:potential-correction-bound}
 0\le q_0(y)\le Ca_y^{-4}\tau_y^{1/2}.
\end{equation}
\end{lemma}

\begin{proof}
The scale function is $1/L$-Lipschitz.  The function
$a\mapsto c_1a\min(a,r_*)^w$ has Lipschitz constant $O(r_*^w)$.
Thus $|t_z-t_y|\le O(r_*^w)|z-y|$.  If $|z-y|\le t_z$, this proves
the support and comparability assertions.  Replacing $t_z$ by $t_y$
changes the kernel by at most $Cr_*^wt_y^{-3}$ on a ball of radius
$2t_y$, giving the first assertion of \eqref{m:eq:master-kernel-bounds}.
The chain rule for a Lipschitz width gives the gradient bounds, first
pointwise almost everywhere and then in $L^2$ by the support bound.
Since $\int\mathcal K_z(y)\,dy=1$, Fubini proves preservation of
integrals.  Weighted Jensen followed by Fubini, using the bound on
$R1$, proves the $L^{5/3}$ estimate.

Newton's theorem gives the nonnegative representation
\[
 q_0(y)=\int\rho_0(z)
   \bigl[K(y-z)-(K*g_{t_z}^2)(y-z)\bigr]\,dz.
\]
Only $|z-y|\le2t_y$ contributes.  The local field bound gives
\[
 \rho_0(z)\le Ca_y^{-6}
    +Ca_y^{-3/2}\sum_{\rm loc}z_m|z-R_m|^{-3/2}.
\]
Equation~\eqref{m:eq:local-power-potential} therefore bounds the displayed
integral by
$C(a_y^{-6}t_y^2+a_y^{-9/2}t_y^{1/2})$, proving
\eqref{m:eq:potential-correction-bound}.  Both $\rho_0$ and $R\rho_0$
belong to $L^1\cap L^{5/3}$, hence also to $L^{6/5}$.  Their potentials
are continuous and decay by the proof of Lemma~\ref{m:lem:barrier-global-tf};
the Sobolev inequality and the Coulomb energy identity put their
difference in $\dot H^1$.
\end{proof}

Define the nonlinearity acting on an offset $h\in\mathbb R$ by
\begin{equation}\label{m:eq:offset-reaction}
 f_y(h)=\cF\int\mathcal K_z(y)(V_c(z)+h)_+^p\,dz.
\end{equation}
For each $y$, this is finite, convex, continuous, and nondecreasing in
$h$.  These properties follow from those of the positive power and the
local integrability of $V_c^p$.  We will use this averaged law in place
of the unsmeared Thomas--Fermi equation.

\subsection{Barriers with a strict gap between scales}

We have established the reference fields and the averaging estimates.
The remaining task is to lower the reference field just enough that its
averaged density pays for the Laplacian of the lowering term.  The
strict inequality between the coefficients below leaves room for all
errors from averaging.

\begin{lemma}[Comparison barriers]\label{m:lem:barrier-subsolution}
There are fixed positive constants $\alpha,\epsilon_0,C'_0,c_b,c'$
with the following properties for all
sufficiently small $r_*>0$.  For $0<r\le r_*$ define
\begin{equation}\label{m:eq:barrier-definition}
 \Gamma_r=\epsilon_0W_s+C'_0r^{4\alpha}W_s^{1+\alpha},\qquad
 B_r=W_0+q_0-\Gamma_r,\qquad b_r=B_r-V_c,
 \qquad h_b=h_0-c_ba^{-4}.
\end{equation}
The offsets $b_r$ are continuous and locally $H^1$ and decay at infinity.
On $\{a\ge r\}$ their Laplacian densities satisfy
\begin{equation}\label{m:eq:barrier-subsolution}
 \Delta b_r\ge-4\pi\nu_s+f_y(b_r).
\end{equation}
More precisely, this inequality holds almost everywhere there between
the locally integrable densities representing the two sides; on every
open set contained in this region it is also a distributional inequality.
On the same region,
\begin{equation}\label{m:eq:barrier-bounds}
 B_r\ge0.9W_0>0,\qquad
 h_b+c'a^{-4}\le b_r\le Ca^{-4}.
\end{equation}
For every contributing center $z$ in \eqref{m:eq:offset-reaction},
\begin{equation}\label{m:eq:reaction-positive-margin}
 V_c(z)+h_b(y)\ge ca_y^{-4}
 \quad\text{almost everywhere in }z.
\end{equation}
All these constants are independent of $r,r_*$ and the nuclear system.
\end{lemma}

\begin{proof}
First fix $\theta$ as in Lemma~\ref{m:lem:barrier-fields}.  That lemma
implies
\[
 D_*:=\sup_x\frac{|\nabla W_s(x)|^2}{W_s(x)^{5/2}}<\infty
\]
with a fixed universal bound.  Since
$\Delta W_s\le \cF W_s^{3/2}$, the chain rule gives
\[
 \Delta(W_s^{1+\alpha})
 \le(1+\alpha)(\cF+\alpha D_*)W_s^{\alpha+3/2}.
\]
Choose $\alpha>0$ so small that
$(1+\alpha)(1+\alpha D_*/\cF)\le1.1$.  Consequently
\begin{equation}\label{m:eq:damping-laplacian}
 \Delta\Gamma_r\le1.1\cF\sqrt{W_s}\,\Gamma_r.
\end{equation}
All chain rules are local: $W_s$ is positive and locally $C^1$ with
locally bounded Laplacian, so multiplication in the weak equation
justifies them.

Choose $c_b>0$ smaller than a sufficiently small fixed fraction of
$c_0$.  On $a\ge r$,
$r^{4\alpha}W_s^\alpha\le C^\alpha$.  We may therefore choose positive
$\epsilon_0,C'_0$ so small that, throughout this region, $\Gamma_r$
is a tiny fixed fraction of $W_s$, is at most $0.1c_0a^{-4}$, and is
at most $(c_b/2)a^{-4}$.  At the same time,
\begin{equation}\label{m:eq:damping-lower}
 \Gamma_r\ge\epsilon_0c_\theta a^{-4}.
\end{equation}
These choices precede every smallness requirement on $r_*$.

The equations for $h_0$ and $q_0$ give the exact cancellation
\begin{equation}\label{m:eq:smoothed-reference-equation}
 \Delta(h_0+q_0)=-4\pi\nu_s+\cF R(W_0^p).
\end{equation}
Fix $y$ with $a_y\ge r$.  For a contributing center $z$, the offset
oscillation estimate and \eqref{m:eq:potential-correction-bound} imply
\begin{align*}
 V_c(z)+b_r(y)
 &=W_0(z)+h_0(y)-h_0(z)+q_0(y)-\Gamma_r(y)\\
 &\le W_0(z)-\Gamma_r(y)+Ca_y^{-4}\tau_y^{1/2}.
\end{align*}
By \eqref{m:eq:damping-lower}, the final error is $o(\Gamma_r(y))$
uniformly as $r_*\to0$.  Furthermore,
\[
 W_0(z)\ge\frac{W_s(z)}{1.01}
      \ge\frac{1-o(1)}{1.01}W_s(y),
\]
using the gradient and lower bounds for $W_s$ and $|z-y|\le2t_y$.
The chosen small ratio $\Gamma_r/W_s$ and the derivative of the power
therefore give
\[
 W_0(z)^p-(V_c(z)+b_r(y))_+^p
       \ge 1.35\sqrt{W_s(y)}\,\Gamma_r(y)
\]
for sufficiently small damping fractions and then sufficiently small
$r_*$.  Indeed the limiting coefficient is
$p/\sqrt{1.01}>1.49$.  After integration in $z$, the factor
$R1(y)=1+o(1)$ leaves the lower bound
$1.3\sqrt{W_s(y)}\Gamma_r(y)$.  Combine this with
\eqref{m:eq:smoothed-reference-equation} and
\eqref{m:eq:damping-laplacian} to prove
\eqref{m:eq:barrier-subsolution}.  All terms represent locally integrable
densities, so the assertion includes points of a level set of $a$ up
to a Lebesgue-null set; no regularity of those level sets is required.

Because $q_0\ge0$ and $\Gamma_r\le0.1W_0$, we have $B_r\ge0.9W_0$.
The upper offset bound follows from $b_r\le h_0+q_0$.  Also
\[
 b_r-h_b=q_0-\Gamma_r+c_ba^{-4}\ge(c_b/2)a^{-4},
\]
so one may take $c'=c_b/2$.  On the kernel support,
\[
 V_c(z)+h_b(y)
   =W_0(z)+h_0(y)-h_0(z)-c_ba_y^{-4}
   \ge ca_y^{-4},
\]
by the lower bound for $W_0$, scale comparability, and the offset
oscillation estimate, after decreasing $r_*$ further.  This proves
\eqref{m:eq:reaction-positive-margin}.

Finally, $h_0,q_0$ are continuous and locally $H^1$, and the same is true
of $\Gamma_r$.  Their decay proves the remaining assertions.
\end{proof}

\Needspace{4\baselineskip}
\begin{lemma}[Gap under doubling]\label{m:lem:barrier-gap}
Fix $\kappa\in(0,1/10)$.  The barriers in
Lemma~\ref{m:lem:barrier-subsolution} have a fixed constant $\gamma>0$ such
that, whenever $R_+=2r\le r_*$,
\begin{equation}\label{m:eq:barrier-gap}
 b_r-b_{R_+}\ge\gamma R_+^{-4}
 \qquad\text{on }\{R_+\le a\le(1+\kappa)R_+\}.
\end{equation}
\end{lemma}

\begin{proof}
The difference is exactly
\[
 b_r-b_{R_+}
   =C'_0(1-2^{-4\alpha})R_+^{4\alpha}W_s^{1+\alpha}.
\]
The lower bound for $W_s$ gives the assertion with
\[
 \gamma=C'_0(1-2^{-4\alpha})
       c_\theta^{1+\alpha}(1+\kappa)^{-4(1+\alpha)}>0.
\]
\end{proof}

The lower margin in \eqref{m:eq:reaction-positive-margin} also controls
inversion of the averaged law.  If $\delta_0>0$ is a sufficiently small
fixed constant, differentiation under the integral gives
\begin{equation}\label{m:eq:barrier-reaction-derivative}
 \frac{d}{dh}\frac{f_y(h)}{4\pi}\ge ca_y^{-2}
 \qquad\text{for }h\ge h_b(y)-\delta_0a_y^{-4}.
\end{equation}
Indeed $(V_c(z)+h)_+^{p-1}\ge ca_y^{-2}$ on the kernel support, and
$R1(y)$ is bounded below.  This derivative bound and the gap
\eqref{m:eq:barrier-gap} are the two margins used in the conditional field
comparisons and the subsequent passage from $r$ to $2r$.
